\begin{textsample}{POS Dim 1 – 2014\_c5 – Score 9.00 – 2014\_c5/enade\_2014\_ci\_bio\_881d1d0b22.txt}  \label{ex:f1_pos_109}
CAPÍTULO 5 DISTRIBUIÇÃO DOS CONCEITOS 5.1 \textbf{PANORAMA} NACIONAL DA DISTRIBUIÇÃO DOS CONCEITOS O Conceito Enade é um indicador de qualidade que avalia o desempenho dos estudantes a partir dos resultados do ENADE.
Os valores possíveis vão de 1 ( pior situação ) a 5 ( melhor situação ).
As Unidades de Observação com apenas um ou sem nenhum concluinte participante não obtêm o Conceito ENADE, ficando Sem Conceito ( SC ).
A Tabela 5.1 apresenta a quantidade e distribuição de cursos de Ciências Biológicas ( Bacharelado ) participantes do ENADE/2014, por faixa de conceito e Grande Região.
Mais uma vez, ressaltando que a diferença entre os cursos \textbf{tabulados} neste capítulo e no capítulo 2 são os cursos sem conceito, em princípio, aqueles sem estudantes concluintes que participassem da prova ou que tivessem tão somente um estudante concluinte22.
Observando -se os dados da Tabela 5.1, nota -se que, dos 227 cursos participantes, 88 ( 38,8 \% ) classificaram -se com conceito 3, o valor modal.
Este foi também o conceito modal em mais três regiões : Nordeste ( 40,0 \% ), Sudeste ( 43,1 \% ) e Sul ( 40,0 \% ).
Nas regiões Norte e Centro-Oeste, a moda foi o conceito 2 ( 46,2 \% e 52,2 \%, respectivamente ).
O conceito 2 foi o segundo mais frequente em nível nacional ( 26,0 \%, correspondendo a 59 cursos ) e o conceito 4, o terceiro ( 22,5 \%, correspondendo a 51 cursos ).
Houve, ainda, 14 cursos ( 6,2 \% ) que receberam conceito 5, e 11 cursos ( 4,8 \% ) que receberam conceito 1.
Dos 227 cursos de Ciências Biológicas ( Bacharelado ), apenas quatro cursos ( 1,8 \% ) ficaram sem conceito ( SC ). 22 Estes cursos com somente um concluinte não foram considerados no capítulo 2.
A região Norte participou com 13 cursos ou 5,7 \% do total nacional.
Nesta região, o conceito 2 foi atribuído a seis dos 13 cursos participantes, o que \textbf{equivale} a 46,2 \% do total regional, além de ser o conceito modal na região.
Nenhum curso ficou sem conceito ( SC ).
Dois cursos foram avaliados com conceito 1 ; três cursos, com conceito 3 ; e dois cursos foram avaliados com conceitos 4 e 5 ( um curso, cada ), correspondendo a 7,7 \% cada.
A região Nordeste participou com 30 cursos ou 13,2 \% do total nacional.
Nesta região, 12 cursos, 40,0 \% em termos regionais, obtiveram conceito 3, o conceito modal para a região.
O conceito 2 foi atribuído a oito cursos ( 26,7 \% ), e ao mesmo número de cursos foi atribuído o conceito 4.
Apenas dois cursos obtiveram conceito 1, correspondendo a 6,7 \% cada, nessa região.
Nenhum curso obteve o conceito 5 ou ficou sem conceito ( SC ).
Dos 116 cursos participantes da região Sudeste ( 51,1 \% do total nacional ), o conceito 3 também foi o conceito modal, tendo sido o obtido por 50 cursos ( 43,1 \% ).
O conceito 1 foi atribuído a dois cursos ( 1,7 \% ) e os conceitos 2 e 4, a 27 e 25 cursos, correspondendo a 23,3 \% e 21,6 \%, respectivamente.
O conceito 5 foi atribuído a nove cursos ( 7,8 \% ).
Três cursos ( 2,6 \% ) na região Sudeste ficaram sem conceito ( SC ).
Nos 45 cursos da região Sul, que correspondem a 19,8 \% do total nacional, há predominância do conceito 3, atribuído a 40,0 \% dos cursos da região ( conceito modal ), ou seja, a 18 dos 45 cursos participantes na região Sul.
O conceito 1 foi atribuído a três cursos ( 6,7 \% ), o conceito 2, a seis cursos ( 13,3 \% ), o conceito 4, a 15 cursos ( 33,3 \% ) e o conceito 5, a três cursos ( 6,7 \% ).
Nenhum curso na região Sul ficou sem conceito ( SC ).
Já dos 23 cursos participantes na região Centro-Oeste, 12 ( 52,2 \% ) receberam conceito 2, o conceito modal.
Apenas um curso ficou sem conceito ( SC ), correspondendo a 4,3 \% dos cursos nessa região.
O conceito 1 foi atribuído a dois cursos ( 8,7 \% ).
A dois outros cursos foi atribuído conceito 4.
Cinco cursos ( 21,7 \% ) obtiveram conceito 3, e apenas um curso ( 4,3 \% ) obteve conceito 5.
O Gráfico 5.1 apresenta a distribuição do Conceito Enade segundo Grande Região : \textbf{poligonais} mais à esquerda representam Grandes Regiões com pior distribuição, e \textbf{poligonais} mais à direita, Grandes Regiões com melhores conceitos.
As regiões Centro-Oeste ( linha preta ) e Norte ( linha azul ) apresentam os cursos em IES com conceitos mais baixos, ao passo que as regiões Sul ( linha \textbf{roxa} ), Sudeste ( linha vermelha ) e Nordeste ( linha verde ) apresentam as distribuições com valores maiores.
A Tabela 5.2 apresenta a quantidade e distribuição de cursos de Ciências Biológicas ( Licenciatura ) participantes do ENADE/2014, por faixa de conceito e Grande Região.
A diferença entre os cursos \textbf{tabulados} neste capítulo e no capítulo 2 são os cursos sem conceito ( SC ), em princípio, aqueles sem estudantes concluintes que participassem da prova ou que tivessem tão somente um estudante concluinte23.
Observando -se os dados da Tabela 5.2, nota -se que, dos 470 cursos participantes, 211 ( 44,9 \% ) classificaram -se com conceito 3, o valor modal, assim como em Ciências Biológicas ( Bacharelado ).
Este foi também o conceito modal em todas as regiões ( entre 42,2 \% na região Nordeste e 56,4 \% na região Centro-Oeste ).
O conceito 2 foi o segundo mais frequente em nível nacional ( 23,6 \%, correspondendo a 111 cursos ) e o conceito 4, o terceiro ( 20,6 \%, correspondendo a 97 cursos ).
Houve, ainda, 19 cursos ( 4,0 \% ) que receberam conceito 5 e outros 21 cursos que receberam conceito 1.
Dos 470 cursos de Ciências Biológicas ( Licenciatura ), 11 cursos ( 2,3 \% ) ficaram sem conceito ( SC ).
A região Norte participou com 53 cursos ou 11,3 \% do total nacional.
Nesta região, o conceito 3 foi atribuído a 23 cursos ( 43,4 \% ) – o mesmo número de cursos a que foi atribuído o conceito 2, o que identifica os conceitos 2 e 3 como os conceitos modais da região.
Nenhum curso ficou sem conceito ( SC ).
Dois cursos ( 3,8 \% ) foram avaliados com conceito 1, e outros cinco ( 9,4 \% ) com conceito 4.
A nenhum curso da região foi atribuído o conceito 5.
A região Nordeste participou com 102 cursos ou 21,7 \% do total nacional.
Nesta região, 43 cursos, 42,2 \% em termos regionais, obtiveram conceito 3, o conceito modal para a região.
Os conceitos 2 e 4 foram atribuídos a, respectivamente, 31 cursos ( 30,4 \% ) e 18 cursos ( 17,6 \% ).
Nove cursos ( 8,8 \% ) obtiveram conceito 1 e 18 cursos ( 17,6 \% ), conceito 4.
Nenhum curso obteve o conceito 5, e somente um curso ( 1,0 \% ) ficou sem conceito ( SC ) nessa região. 23 Estes cursos com somente um concluinte não foram considerados no capítulo 2.
Dos 180 cursos participantes da região Sudeste ( 38,3 \% do total nacional ), o conceito 3, conceito modal, foi obtido por 78 cursos ( 43,3 \% ).
O conceito 1 foi atribuído a cinco cursos ( 2,8 \% ), e os conceitos 2 e 4, respectivamente, a 35 cursos ( 19,4 \% ) e a 45 cursos ( 25,0 \% ).
O conceito 5 foi atribuído a 11 cursos ( 6,1 \% ).
Seis cursos ( 3,3 \% ) ficaram sem conceito ( SC ) nessa região.
Os 80 cursos da região Sul ficaram distribuídos em todas as faixas de conceitos.
A predominância foi do conceito 3, atribuído a 45,0 \% dos cursos ( conceito modal ), ou seja, a 36 dos 80 cursos participantes na região Sul.
O conceito 1 foi atribuído a dois cursos ( 2,5 \% ), o conceito 2, a dez cursos ( 12,5 \% ), o conceito 4, a 20 cursos ( 25,0 \% ) e o conceito 5, a oito cursos ( 10,0 \% ).
Quatro cursos ( 5,0 \% ) na região Sul ficaram sem conceito ( SC ).
Já dos 55 cursos participantes na região Centro-Oeste, 31 ( 56,4 \% ) receberam conceito 3, o conceito modal.
Nesta região, nenhum curso ficou sem conceito ( SC ) ou com conceito 5.
A 12 cursos ( 21,8 \% ) foi atribuído conceito 2, enquanto nove cursos ( 16,4 \% ) obtiveram conceito 4, nessa região.
Três cursos ( 5,5 \% ) foram avaliados com conceito 1.
O Gráfico 5.2 apresenta a distribuição do Conceito Enade segundo Grande Região : \textbf{poligonais} mais à esquerda representam Grandes Regiões com pior distribuição, e \textbf{poligonais} mais à direita, Grandes Regiões com melhores conceitos.
As regiões Norte ( linha azul ) e Nordeste ( linha verde ) apresentam os cursos em IES com conceitos mais baixos, ao passo que as regiões Sudeste ( linha vermelha ) e Sul ( linha \textbf{roxa} ) apresentam as distribuições com valores maiores.
A região Centro-Oeste ( linha preta ) apresenta uma situação intermediária. 5.2 CONCEITOS POR CATEGORIA ADMINISTRATIVA E POR GRANDE REGIÃO A Tabela 5.3 apresenta a distribuição dos cursos participantes do ENADE/2014 de Ciências Biológicas ( Bacharelado ), por Categoria Administrativa, de acordo com os conceitos por eles alcançados, segundo as Grandes Regiões brasileiras.
Dos 227 cursos participantes, 93 ( 41,0 \% ) eram \textbf{ministrados} em Instituições \textbf{Públicas} e 134 ( 59,0 \% ), em Instituições Privadas.
De acordo com as informações da Tabela 5.3, em termos nacionais, entre Instituições \textbf{Públicas} e Privadas, das 14 IES cujos cursos receberam conceito 5, 12 eram Instituições \textbf{Públicas} e duas, Instituições Privadas.
Dos 93 cursos participantes de IES Públicas, o conceito 4 foi o valor modal, atribuído a 35 cursos ( 37,6 \% ).
Entre os demais cursos de Instituições \textbf{Públicas} participantes, cinco cursos ( 5,4 \% ) foram avaliados com conceito 1, onze cursos ( 11,8 \% ) obtiveram conceito 2, 30 cursos ( 32,3 \% ) receberam conceito 3 e a doze cursos ( 12,9 \% ) foi atribuído o conceito 5.
Nesta categoria, nenhum curso ficou sem conceito ( SC ).
Na rede privada, o conceito modal foi 3, com 58 cursos ( 43,3 \% ) dos 134 cursos da categoria.
Entre os demais cursos participantes, seis ( 4,5 \% ) receberam conceito 1, e 48 ( 35,8 \% ), conceito 2.
O conceito 4 foi atribuído a 16 cursos ( 11,9 \% ) e o conceito 5, a dois cursos ( 1,5 \% ).
Nesta Categoria Administrativa, outros quatro cursos ficaram sem conceito ( SC ).
Na análise por região, observa -se que, na região Norte, as Instituições Privadas participaram com seis cursos ( 46,2 \% do total regional ), dos quais nenhum ficou sem conceito ( SC ) ou com conceitos 4 ou 5.
O conceito modal para as Instituições Privadas na região foi 2, correspondendo a quatro cursos ( 66,7 \% ).
Dos dois cursos \textbf{restantes}, um obteve conceito 1, e o outro, conceito 3.
As Instituições Públicas participaram com sete cursos na região Norte ( 53,8 \% do total regional ), dos quais dois obtiveram conceito 2, e outros dois, o conceito 3, valores modais da distribuição.
Os conceitos 1, 4 e 5 foram atribuídos a cada um dos três cursos \textbf{restantes}, correspondendo, cada qual, a 14,3 \% do total da região.
Nenhum curso ficou sem conceito ( SC ).
Na região Nordeste, a rede privada concentrou 5 dos 30 cursos participantes, o equivalente a 16,7 \% do total da região, a menor das proporções desta categoria considerando todas as regiões do país.
Dentre os cursos oferecidos por IES Privadas no Nordeste, um curso ficou alocado ao conceito 3.
Os quatro \textbf{restantes} obtiveram conceito 2, não ocorrendo casos de cursos que tenham ficado sem conceito ( SC ).
As Instituições \textbf{Públicas} dessa região participaram com 25 cursos ( 83,3 \% ), dos quais 11 ( 44,0 \% ) obtiveram conceito 3, o valor modal, seguidos por oito cursos que receberam conceito 4 ( 32,0 \% ).
Dois cursos foram alocados ao conceito 1, enquanto outros quatro obtiveram conceito 2.
Nenhum curso recebeu conceito 5 ou ficou sem conceito ( SC ) nesta combinação de categoria e região.
Na região Sudeste, a proporção de cursos da rede privada foi 69,0 \%, correspondendo a 80 dos 116 cursos participantes.
Nesta categoria, na região Sudeste, o conceito modal foi 3 ( 40 cursos ).
Nesta combinação de Categoria Administrativa e Grande Região, três cursos não receberam conceito ( SC ).
Os demais foram avaliados com conceito 1 ( dois cursos ), conceito 2 ( 25 cursos ), conceito 4 ( oito cursos ) e conceito 5 ( dois cursos ).
Entre os 36 cursos oferecidos por Instituições \textbf{Públicas} na região Sudeste, o conceito modal foi 4 ( 17 cursos ), e os demais cursos receberam os conceitos 2 ( dois cursos ), 3 ( dez cursos ) e 5 ( sete cursos ).
Na região Sudeste, nesta categoria, nenhum curso recebeu conceito 1 ou ficou sem conceito ( SC ).
As Instituições Privadas concentraram 27 dos 45 cursos participantes da região Sul, 60,0 \% do total regional.
Desses, 13 ficaram com conceito 3, o conceito modal.
Nesta combinação de Categoria Administrativa e Grande Região, seis cursos ficaram com conceito 2, e outros seis, com conceito 4.
Dois cursos obtiveram conceito 1, enquanto nenhum curso obteve conceito 5 ou ficou sem conceito ( SC ).
As Instituições \textbf{Públicas} na região Sul participaram com 18 cursos ( 40,0 \% ) e o conceito modal foi 4, atribuído a nove cursos.
Um único curso foi avaliado com conceito 1, cinco cursos receberam conceito 3 e três cursos obtiveram conceito 5.
Nenhum curso ficou sem conceito ( SC ) ou recebeu conceito 2.
Na região Centro-Oeste, 16 dos 23 cursos participantes eram de Instituições Privadas ( 69,6 \% em termos regionais ), a proporção mais elevada entre as regiões brasileiras.
Destes, nove concentraram -se no conceito 2, conceito modal.
Dentre os demais, um curso recebeu conceito 1, três cursos, conceito 3, e dois cursos, conceito 4.
Nessa categoria, um curso ficou sem conceito ( SC ) e nenhum recebeu conceito 5.
Dos sete cursos oferecidos por Instituições \textbf{Públicas}, três foram avaliados com conceito 2, conceito modal, seguidos de dois cursos que obtiveram conceito 3.
Um único curso recebeu conceito 5.
Também a um único curso foi atribuído conceito 1.
Nesta região, nenhum curso de IES Pública ficou sem conceito ( SC ) ou recebeu conceito 4.
O Gráfico 5.3 apresenta a distribuição do Conceito Enade segundo a Categoria Administrativa da IES.
Os cursos em IES Públicas ( linha azul ) apresentam uma \textbf{poligonal} mais à direita do que os em IES Privadas ( linha verde ) e, consequentemente, uma distribuição de Conceitos Enade com valores maiores.
A Tabela 5.4 apresenta a distribuição dos cursos participantes do ENADE/2014 de Ciências Biológicas ( Licenciatura ), por Categoria Administrativa, de acordo com os conceitos por eles alcançados, segundo as Grandes Regiões brasileiras.
Dos 470 cursos participantes, 241 ( 51,3 \% ) eram \textbf{ministrados} em Instituições \textbf{Públicas} e 229 ( 48,7 \% ), em Instituições Privadas.
De acordo com as informações da Tabela 5.4, em termos nacionais, entre Instituições \textbf{Públicas} e Privadas, das 19 IES que receberam conceito 5, 15 eram Públicas e quatro, Privadas.
Dos 241 cursos participantes de IES Públicas, o conceito 3 foi o valor modal, atribuído a 104 cursos ( 43,2 \% ).
Entre os demais cursos de Instituições \textbf{Públicas} participantes, nenhum curso ficou sem conceito ( SC ), 62 cursos receberam conceito 4 ( 25,7 \% ) e, como já comentado, 15 ( 6,2 \% ) foram alocados ao conceito 5.
Nesta categoria, dez cursos ( 4,1 \% ) obtiveram conceito 1, enquanto 50 cursos ( 20,7 \% ) obtiveram conceito 2.
Na rede privada, o conceito modal também foi 3, com 107 cursos dos 229 da categoria.
Entre os cursos participantes, 11 receberam conceito 1 e 61, conceito 2.
O conceito 4 foi atribuído a 35 cursos e o conceito 5, a quatro cursos.
Nesta Categoria Administrativa, 11 cursos ficaram sem conceito ( SC ).
Na análise por região, observa -se que, na região Norte, as Instituições Privadas participaram com 19 cursos ( 35,8 \% do total regional ), dos quais nenhum ficou sem conceito ( SC ) ou com conceito 5.
O conceito modal para as Instituições Privadas na região foi 2, com 12 cursos, correspondendo a 63,2 \%.
O conceito 4 foi atribuído a um único curso, e o conceito 3, a cinco cursos.
Também a um único curso foi atribuído o conceito 1 nessa combinação de Categoria Administrativa e Grande Região.
As Instituições Públicas participaram com 34 cursos na região Norte ( 64,2 \% do total regional ).
Desses, 18 cursos obtiveram conceito 3, conceito modal nessa região para esta categoria.
Quatro cursos obtiveram conceito 4, onze cursos receberam conceito 2 e a apenas um curso foi atribuído conceito 1.
Nessa categoria, nenhum curso obteve conceito 5 ou ficou sem conceito ( SC ). atribuído a um único curso, e o conceito 3, a cinco cursos.
Também a um único curso foi atribuído o conceito 1 nessa combinação de Categoria Administrativa e Grande Região.
As Instituições Públicas participaram com 34 cursos na região Norte ( 64,2 \% do total regional ).
Desses, 18 cursos obtiveram conceito 3, conceito modal nessa região para esta categoria.
Quatro cursos obtiveram conceito 4, onze cursos receberam conceito 2 e a apenas um curso foi atribuído conceito 1.
Nessa categoria, nenhum curso obteve conceito 5 ou ficou sem conceito ( SC ).
Na região Nordeste, a rede privada concentrou 16 dos 102 cursos participantes, o equivalente a 15,7 \% do total da região, a menor proporção desta categoria dentre todas as regiões do país.
Entre os cursos oferecidos por IES Privadas no Nordeste, predominaram os cursos com conceito 3 : seis cursos, correspondendo a 37,5 \% ( conceito modal ).
Quatro cursos ( 25,0 \% ) receberam conceito 1, seguidos por mais quatro cursos ( 25,0 \% ) que receberam conceito 2, enquanto o conceito 4 foi atribuído a somente um curso ( 6,3 \% ).
Nenhum curso obteve conceito 5, enquanto um curso ficou sem conceito ( SC ).
As Instituições \textbf{Públicas} dessa região participaram com 86 cursos ( 84,3 \% ), a maior proporção desta categoria dentre todas as regiões do país, dos quais 37 cursos ( 43,0 \% ) obtiveram conceitos 3, o valor modal nessa combinação de Categoria Administrativa e Grande Região.
Cinco cursos ( 5,8 \% ) receberam conceito 1, seguidos de 27 cursos ( 31,4 \% ) que obtiveram conceito 2.
O conceito 4 foi atribuído a 17 cursos ( 19,8 \% ), não se registrando casos de cursos que tenham recebido conceito 5 ou que tenham ficado sem conceito ( SC ) nesta combinação de categoria e região.
Na região Sudeste, a proporção de cursos da rede privada, 66,7 \%, foi a mais elevada considerando todas regiões brasileiras, correspondendo a 120 dos 180 cursos participantes.
Nesta categoria, na região Sudeste, o conceito modal foi 3 ( 59 cursos ).
Nesta combinação de Categoria Administrativa e Grande Região, seis cursos ficaram sem conceito ( SC ).
Os demais foram avaliados com conceito 1 ( três cursos ), conceito 2 ( 30 cursos ), conceito 4 ( 20 cursos ) e conceito 5 ( dois cursos ).
Entre os 60 cursos oferecidos por Instituições \textbf{Públicas} na região Sudeste, o conceito modal foi 4 ( 25 cursos ).
Os demais cursos receberam os conceitos 1 ( dois cursos ), 2 ( cinco cursos ), 3 ( 19 cursos ) e 5 ( nove cursos ).
Na região Sudeste, nesta categoria, nenhum curso ficou sem conceito ( SC ).
As Instituições Privadas concentraram 53 dos 80 cursos participantes da região Sul, 66,3 \% do total regional, a segunda proporção mais elevada registrada entre as regiões do país.
Desses, 26 cursos ficaram com conceito 1, nove cursos obtiveram conceito 2, dez cursos receberam conceito 4 e, por fim, dois cursos foram avaliados com conceito 5.
Quatro cursos ficaram sem conceito ( SC ).
As Instituições \textbf{Públicas} na região Sul participaram com 27 cursos, 33,8 \% do total da região.
Os conceitos modais foram 3 e 4, com dez cursos cada.
Os demais cursos foram avaliados com conceitos 2 ( um curso ) e 5 ( seis cursos ).
Nessa categoria, nenhum curso recebeu conceito 1 ou ficou sem conceito ( SC ).
Na região Centro-Oeste, 21 dos 55 cursos participantes eram de Instituições Privadas ( 38,2 \% em termos regionais ).
Destes, 11 concentraram -se no conceito 3, conceito modal.
Dentre os demais, somente um curso recebeu conceito 1, seis cursos receberam conceito 2, e três cursos receberam conceito 4.
Nenhum curso recebeu conceito 5 ou ficou sem conceito ( SC ) nesta categoria.
Dos 34 cursos oferecidos por Instituições \textbf{Públicas}, 20 foram avaliados com conceito 3, conceito modal.
Seis cursos obtiveram conceito 2, e outros seis, conceito 4.
Dois cursos foram avaliados com conceito 1, e a nenhum curso foi atribuído conceito 5.
Nesta região, nenhum curso de IES \textbf{pública} ficou sem conceito ( SC ).
O Gráfico 5.4 apresenta a distribuição do Conceito Enade segundo a Categoria Administrativa da IES.
Os cursos em IES Públicas ( linha azul ) apresentam uma \textbf{poligonal} mais à direita do que os cursos em IES Privadas ( linha verde ) e, consequentemente, uma distribuição de Conceitos Enade com valores maiores. 5.3 CONCEITOS POR ORGANIZAÇÃO ACADÊMICA E POR GRANDE REGIÃO Na Tabela 5.5, encontra -se a distribuição dos conceitos atribuídos aos cursos participantes do ENADE/2014 na Área de Ciências Biológicas ( Bacharelado ), por Organização Acadêmica, segundo as Grandes Regiões brasileiras.
Dos 227 cursos de Ciências Biológicas ( Bacharelado ) participantes, 158 eram oferecidos em Universidades, 37 cursos em Centros Universitários e 32 em Faculdades.
Esta distribuição corresponde a, respectivamente, 69,6 \%, 16,3 \% e 14,1 \% dos cursos.
De acordo com os dados apresentados, de todos os 14 cursos avaliados com conceito 5, 13 eram vinculados a Universidades.
Apenas dois cursos oferecidos em Universidades ficaram sem conceito ( SC ).
Esse tipo de Organização Acadêmica teve o conceito 3 como modal, com 57 cursos ( 36,1 \% ).
Os demais cursos de Universidades avaliados receberam os conceitos 1 ( nove cursos ), 2 ( 29 cursos ), 4 ( 48 cursos ) e conceito 5 ( 13 cursos, como já mencionado ).
Entre os cursos em Centros Universitários, o conceito modal foi 3, com 24 cursos ( 64,9 \% ).
Neste tipo de Organização Acadêmica, dois cursos ficaram sem conceito ( SC ), e os outros cursos receberam os conceitos 2 ( nove cursos ) e 4 ( dois cursos ).
Nenhum curso recebeu conceito 1 ou 5.
Nas Faculdades, nenhum dos 32 cursos ficou sem conceito ( SC ).
Dois cursos ficaram alocados no conceito 1, seguido de 21 cursos ( 65,6 \% ) que obtiveram conceito 2, conceito modal, e de sete cursos que receberam conceito 3.
Os demais cursos neste tipo de Organização Acadêmica receberam os conceitos 4 e 5 ( um curso cada ).
Considerando -se separadamente as regiões brasileiras, verifica -se que, na região Norte, as Universidades concentraram sete dos 13 cursos participantes.
Os conceitos modais foram os conceitos 2 e 3, com dois cursos cada.
Os demais conceitos ( 1, 4 e 5 ) foram atribuídos, cada qual, a um curso dos três cursos \textbf{restantes}.
Nenhum curso ficou sem conceito ( SC ) nesta categoria e região.
Os Centros Universitários da região Norte foram representados por apenas dois cursos, que receberam conceito 2.
As Faculdades participaram com quatro cursos na região Norte, os quais receberam conceitos 1 ( um curso ), 2 ( dois cursos, conceito modal ) e 3 ( um curso ).
Na região Nordeste, as Universidades participaram com 27 dos 30 cursos na Área de Ciências Biológicas ( Bacharelado ).
O conceito modal foi 3, com 12 cursos, enquanto o conceito 4 foi atribuído a oito cursos.
Cinco cursos obtiveram conceito 2, e dois cursos foram alocados no conceito 1.
Nenhum dos cursos oferecidos por Universidades no Nordeste recebeu conceito 5 ou ficou sem conceito ( SC ).
Os Centros Universitários \textbf{contaram} com somente um curso participante na região Nordeste, o qual obteve conceito 2.
As Faculdades foram representadas por dois cursos na região Nordeste, ambos avaliados com conceito 2.
Na região Sudeste, as Universidades concentraram 73 dos 116 cursos da região.
Entre os cursos em Universidades, nessa região, o conceito modal foi 3 com 26 cursos.
Os demais cursos receberam os conceitos 1 ( dois cursos ), 2 ( 13 cursos ), 4 ( 23 cursos ) e 5 ( oito cursos ).
Apenas um curso ficou sem conceito ( SC ).
Os Centros Universitários participaram com 26 cursos na região Sudeste, dos quais 19 obtiveram o conceito modal, 3.
Nenhum curso recebeu conceitos 1 ou 5 nessa categoria e região.
Os demais cursos receberam os conceitos 2 ( quatro cursos ) e 4 ( um curso ).
Dois cursos ficaram sem conceito ( SC ).
As Faculdades foram representadas por 17 cursos na região Sudeste dos quais nenhum recebeu conceito 1, dez receberam conceito 2 ( conceito modal ) e cinco receberam conceito 3.
Os demais cursos se distribuíram nos conceitos 4 e 5 ( um curso cada ).
Nenhum curso ficou sem conceito ( SC ) nessa categoria e região.
Dos 45 cursos da região Sul, 38 eram de Universidades, para os quais o conceito modal foi 4, com 15 cursos.
Os demais cursos receberam os conceitos 1 ( três cursos ), 2 ( três cursos ), 3 ( 14 cursos ) e 5 ( três cursos ).
Nesse tipo de organização, nenhum dos cursos ficou sem conceito ( SC ).
Os três cursos participantes de Centros Universitários da região Sul obtiveram conceito 3.
Não houve casos de cursos sem conceito ( SC ) nessa categoria e região.
Já as Faculdades foram representadas na região Sul por quatro cursos, três dos quais receberam conceito 2, o conceito modal.
Um curso recebeu conceito 3, enquanto nenhum ficou sem conceito ( SC ) ou recebeu os demais conceitos nessa categoria e região.
Na região Centro-Oeste, 13 dos 23 cursos eram de Universidades.
Nesse tipo de organização, somente um curso ficou sem conceito ( SC ), e o conceito modal foi 2, com 6 cursos.
Três cursos obtiveram conceito 3, e os três cursos \textbf{restantes} obtiveram os conceitos 1, 4 e 5 ( um curso cada ).
Os Centros Universitários da região Centro-Oeste \textbf{contaram} com cinco cursos, que obtiveram conceitos 2 ( dois cursos ), 3 ( dois cursos ) e 4 ( um curso ).
Nenhum curso ficou sem conceito ( SC ) ou recebeu os conceitos 1 ou 5.
Dos cinco cursos oferecidos em Faculdades na região Centro-Oeste, quatro foram avaliados com conceito 2 ( conceito modal ), e um com conceito 3.
O Gráfico 5.5 apresenta a distribuição do Conceito Enade segundo a Organização Acadêmica da IES.
Os cursos em Universidades ( linha azul ) apresentam uma \textbf{poligonal} mais à direita do que as demais, \textbf{denotando} uma melhor distribuição de Conceitos ENADE.
No outro extremo, temos os cursos oferecidos em Faculdades ( linha vermelha ) com a \textbf{poligonal} mais à esquerda e, consequentemente, uma distribuição de Conceitos Enade com menores valores.
IES de Centros Universitários ( linha verde ) apresentam uma situação intermediária.
Na Tabela 5.6, encontra -se a distribuição dos conceitos atribuídos aos cursos participantes do ENADE/2014 na Área de Ciências Biológicas ( Licenciatura ), por Organização Acadêmica, segundo as Grandes Regiões brasileiras.
Dos 470 cursos de Ciências Biológicas ( Licenciatura ) participantes, 316 eram oferecidos em Universidades, 102 em Faculdades e os demais 52, em Centros Universitários.
Esta distribuição corresponde a, respectivamente, 67,2 \%, 21,7 \% e 11,1 \% dos cursos.
De acordo com os dados apresentados, 17 dos 19 cursos avaliados com conceito 5 eram vinculados a Universidades.
Seis cursos oferecidos em Universidades ficaram sem conceito ( SC ).
Esse tipo de Organização Acadêmica teve o conceito 3 como modal, com 142 cursos ( 44,9 \% ).
Os demais cursos de Universidades avaliados receberam os conceitos 1 ( dez cursos ), 2 ( 57 cursos ), 4 ( 84 cursos ) e conceito 5 ( 17 cursos, como já mencionado ).
Entre os cursos em Centros Universitários, o conceito modal foi 3, com 29 cursos ( 55,8 \% ).
Neste tipo de Organização Acadêmica, os conceitos 2 e 4 foram atribuídos, respectivamente, a onze e a nove cursos.
A cada um dos dois cursos \textbf{restantes} foi atribuído o conceito 1 e o 5.
Apenas um curso ficou sem conceito ( SC ).
Nas Faculdades, dos 102 cursos participantes, quatro ficaram sem conceito ( SC ).
O conceito modal foi o conceito 2, atribuído a 43 cursos, seguidos por outros 40 cursos que obtiveram conceito 3.
Entre os demais cursos vinculados a esse tipo de Organização Acadêmica, dez receberam conceito 1, quatro receberam conceito 4, e um recebeu conceito 5.
Considerando -se separadamente as regiões brasileiras, verifica -se que, na região Norte, as Universidades concentraram 35 dos 53 cursos participantes.
Dezoito cursos receberam conceito 3 ( conceito modal ), seguidos por 12 cursos que receberam conceito 2.
O conceito 1 foi atribuído a apenas um curso, e o conceito 4, a quatro.
Nenhum curso ficou sem conceito ( SC ) ou recebeu conceito 5 nesta categoria e região.
Os Centros Universitários da região Norte foram representados por quatro cursos.
Destes, dois receberam conceito 3, e outros dois, conceito 2.
Nenhum curso ficou sem conceito ( SC ) ou recebeu conceitos 1, 4 ou 5.
As Faculdades participaram com 14 cursos na região Norte.
Os cursos oferecidos por Faculdades receberam os conceitos 1 ( um curso ), 2 ( nove cursos, conceito modal ), 3 ( três cursos ) e conceito 4 ( um curso ).
Nenhum curso ficou sem conceito ( SC ) ou recebeu conceito 5.
Na região Nordeste, as Universidades participaram com 83 dos 102 cursos na Área de Ciências Biológicas ( Licenciatura ).
O conceito modal foi 3, atribuído a 39 cursos, enquanto os demais cursos oferecidos por Universidades receberam conceitos 1 ( quatro cursos ), 2 ( 23 cursos ) e 4 ( 17 cursos ).
Nenhum curso ficou sem conceito ( SC ) ou recebeu conceito 5.
Os Centros Universitários \textbf{contaram} com três cursos participantes na região Nordeste, os quais ficaram distribuídos igualmente nos conceitos 2, 3 e 4.
Já as Faculdades \textbf{contaram} com 16 cursos participantes nessa região, os quais receberam conceitos 1 ( cinco cursos ), 2 ( sete cursos, conceito modal ) e 3 ( três cursos ).
Um curso ficou sem conceito ( SC ) e nenhum curso recebeu conceitos 4 ou 5.
Na região Sudeste, as Universidades concentraram 99 dos 180 cursos da região.
Entre os cursos em Universidades, na região, o conceito modal foi 4, atribuído a 38 cursos, seguidos de outros 36 cursos que receberam conceito 3.
Três cursos ficaram sem conceito ( SC ) e os demais receberam os conceitos 1 ( um curso ), 2 ( 11 cursos ) e 5 ( dez cursos ).
Os Centros Universitários participaram com 31 cursos na região Sudeste, dos quais 20 obtiveram conceito modal, 3.
Apenas um curso ficou sem conceito ( SC ) e nenhum recebeu os conceitos 1 ou 5 nessa categoria e região.
Os demais receberam os conceitos 2 ( seis cursos ) e 4 ( quatro cursos ).
As Faculdades foram representadas por 50 cursos na região Sudeste, dos quais 22 receberam conceito 3, conceito modal.
Os demais se distribuíram nos conceitos 1 ( quatro cursos ), 2 ( 18 cursos ), 4 ( três cursos ) e 5 ( um curso ).
Dois cursos ficaram sem conceito ( SC ) nessa categoria e região.
Dos 80 cursos da região Sul, 59 eram de Universidades, para os quais o conceito modal foi 3, atribuído a 24 cursos.
Nesse tipo de organização, três cursos ficaram sem conceito ( SC ).
Os demais receberam os conceitos 1 ( dois cursos ), 2 ( cinco cursos ), 4 ( 18 cursos ) e 5 ( sete cursos ).
Os Centros Universitários da região Sul tiveram três dos sete cursos participantes avaliados com o conceito modal, 3.
Nenhum curso ficou sem conceito ( SC ) ou recebeu conceito 1.
Os demais cursos receberam os conceitos 2 ( um curso ), 4 ( dois cursos ) e 5 ( um curso ).
Foram 14 os cursos vinculados a Faculdades na região Sul, dos quais nove receberam conceito 3, o conceito modal, e quatro receberam conceito 2.
Apenas um curso ficou sem conceito ( SC ) e nenhum curso recebeu os conceitos 1, 4 ou 5 nessa categoria e região.
Na região Centro-Oeste, 40 dos 55 cursos eram de Universidades.
Nesse tipo de organização, nenhum curso ficou sem conceito ( SC ) ou com conceito 5.
O conceito modal foi 3, atribuído a 25 cursos.
Os demais receberam os conceitos 1 ( dois cursos ), 2 ( seis cursos ) e 4 ( sete cursos ).
Os Centros Universitários da região Centro-Oeste \textbf{contaram} com sete cursos, dos quais três cursos obtiveram conceito 3, conceito modal.
Os demais cursos receberam os conceitos 1 ( um curso ), 2 ( um curso ) e 4 ( dois cursos ).
Nenhum curso ficou sem conceito ( SC ) ou recebeu conceito 5.
Dos oito cursos oferecidos em Faculdades na região Centro-Oeste, cinco ficaram alocados no conceito 2, conceito modal, e três cursos receberam conceito 3.
Nenhum curso ficou sem conceito ( SC ) ou recebeu os conceitos 1, 4 ou 5 nessa categoria e região.
O Gráfico 5.6 apresenta a distribuição do Conceito Enade segundo a Organização Acadêmica da IES.
Os cursos em Universidades ( linha azul ) apresentam uma \textbf{poligonal} mais à direita do que as demais, \textbf{denotando} uma melhor distribuição de Conceitos ENADE.
No outro extremo, temos os cursos oferecidos em Faculdades ( linha vermelha ) com a \textbf{poligonal} mais à esquerda e, consequentemente, uma distribuição de Conceitos Enade com menores valores.
IES de Centros Universitários ( linha verde ) apresentam uma situação intermediária.
Considerando -se separadamente as regiões brasileiras, verifica -se que, na região Norte, as Universidades concentraram 35 dos 53 cursos participantes.
Dezoito cursos receberam conceito 3 ( conceito modal ), seguidos por 12 cursos que receberam conceito 2.
O conceito 1 foi atribuído a apenas um curso, e o conceito 4, a quatro.
Nenhum curso ficou sem conceito ( SC ) ou recebeu conceito 5 nesta categoria e região.
Os Centros Universitários da região Norte foram representados por quatro cursos.
Destes, dois receberam conceito 3, e outros dois, conceito 2.
Nenhum curso ficou sem conceito ( SC ) ou recebeu conceitos 1, 4 ou 5.
As Faculdades participaram com 14 cursos na região Norte.
Os cursos oferecidos por Faculdades receberam os conceitos 1 ( um curso ), 2 ( nove cursos, conceito modal ), 3 ( três cursos ) e conceito 4 ( um curso ).
Nenhum curso ficou sem conceito ( SC ) ou recebeu conceito 5.
Na região Nordeste, as Universidades participaram com 83 dos 102 cursos na Área de Ciências Biológicas ( Licenciatura ).
O conceito modal foi 3, atribuído a 39 cursos, enquanto os demais cursos oferecidos por Universidades receberam conceitos 1 ( quatro cursos ), 2 ( 23 cursos ) e 4 ( 17 cursos ).
Nenhum curso ficou sem conceito ( SC ) ou recebeu conceito 5.
Os Centros Universitários \textbf{contaram} com três cursos participantes na região Nordeste, os quais ficaram distribuídos igualmente nos conceitos 2, 3 e 4.
Já as Faculdades \textbf{contaram} com 16 cursos participantes nessa região, os quais receberam conceitos 1 ( cinco cursos ), 2 ( sete cursos, conceito modal ) e 3 ( três cursos ).
Um curso ficou sem conceito ( SC ) e nenhum curso recebeu conceitos 4 ou 5.
Na região Sudeste, as Universidades concentraram 99 dos 180 cursos da região.
Entre os cursos em Universidades, na região, o conceito modal foi 4, atribuído a 37 cursos, seguidos de outros 36 cursos que receberam conceito 2.
Três cursos ficaram sem conceito ( SC ) e os demais receberam os conceitos 1 ( um curso ), 2 ( 12 cursos ) e 5 ( dez cursos ).
Os Centros Universitários participaram com 31 cursos na região Sudeste, dos quais 20 obtiveram conceito modal, 3.
Apenas um curso ficou sem conceito ( SC ) e nenhum recebeu os conceitos 1 ou 5 nessa categoria e região.
Os demais receberam os conceitos 2 ( seis cursos ) e 4 ( quatro cursos ).
As Faculdades foram representadas por 46 cursos na região Sudeste, dos quais 22 receberam conceito 3, conceito modal.
Os demais se distribuíram nos conceitos 1 ( cinco cursos ), 2 ( 17 cursos ), 4 ( quatro cursos ) e 5 ( um curso ).
Somente um curso ficou sem conceito ( SC ) nessa categoria e região.
Dos 80 cursos da região Sul, 59 eram de Universidades, para os quais o conceito modal foi 3, atribuído a 24 cursos.
Nesse tipo de organização, três cursos ficaram sem conceito ( SC ).
Os demais receberam os conceitos 1 ( dois cursos ), 2 ( cinco cursos ), 4 ( 18 cursos ) e 5 ( sete cursos ).
Os Centros Universitários da região Sul tiveram três dos sete cursos participantes avaliados com o conceito modal, 3.
Nenhum curso ficou sem conceito ( SC ) ou recebeu conceito 1.
Os demais cursos receberam os conceitos 2 ( um curso ), 4 ( dois cursos ) e 5 ( um curso ).
Foram 14 os cursos vinculados a Faculdades na região Sul, dos quais nove receberam conceito 3, o conceito modal, e quatro receberam conceito 2.
Apenas um curso ficou sem conceito ( SC ) e nenhum curso recebeu os conceitos 1, 4 ou 5 nessa categoria e região.
Na região Centro-Oeste, 40 dos 55 cursos eram de Universidades.
Nesse tipo de organização, nenhum curso ficou sem conceito ( SC ) ou com conceito 5.
O conceito modal foi 3, atribuído a 25 cursos.
Os demais receberam os conceitos 1 ( dois cursos ), 2 ( seis cursos ) e 4 ( sete cursos ).
Os Centros Universitários da região Centro-Oeste \textbf{contaram} com sete cursos, dos quais três cursos obtiveram conceito 3, conceito modal.
Os demais cursos receberam os conceitos 1 ( um curso ), 2 ( um curso ) e 4 ( dois cursos ).
Nenhum curso ficou sem conceito ( SC ) ou recebeu conceito 5.
Dos oito cursos oferecidos em Faculdades na região Centro-Oeste, cinco ficaram alocados no conceito 2, conceito modal, e três cursos receberam conceito 3.
Nenhum curso ficou sem conceito ( SC ) ou recebeu os conceitos 1, 4 ou 5 nessa categoria e região.
O Gráfico 5.6 apresenta a distribuição do Conceito Enade segundo a Organização Acadêmica da IES.
Os cursos em Universidades ( linha azul ) apresentam uma \textbf{poligonal} mais à direita do que as demais, \textbf{denotando} uma melhor distribuição de Conceitos ENADE.
No outro extremo, temos os cursos oferecidos em Faculdades ( linha vermelha ) com a \textbf{poligonal} mais à esquerda e, consequentemente, uma distribuição de Conceitos Enade com menores valores.
IES de Centros Universitários ( linha verde ) apresentam uma situação intermediária.

% matched lemmas: contar, denotar, equivaler, ministrar, panorama, poligonal, pública, restante, roxa, tabular
\end{textsample}
